\editionPage[\pdfbookmark[0]{Indications des défis}{bm-indices-debut-parent}]{indices-debut}{Indications progressives}{Défis 1 à 3}{Ne lire que l’aide nécessaire.}
\begin{blocEdition}{indices-debut}
\editionEtape{D01 · Une commande au meilleur prix}\textbf{Indice 1.} Combien de lots de cinq faut-il pour en avoir au moins 18 ?\par\medskip\textbf{Indice 2.} Comparer les deux dépenses totales, en ajoutant les frais une seule fois.\par
\editionEtape{D02 · Deux restes et une divisibilité}\textbf{Indice 1.} Lister les multiples de 9 dans l’intervalle.\par\medskip\textbf{Indice 2.} Tester le reste de la division par 7 pour chacun.\par
\editionEtape{D03 · Un nombre et son opposé}\textbf{Indice 1.} Si l’opposé est strictement négatif, de quel signe est $a$ ?\par\medskip\textbf{Indice 2.} Pour un nombre positif, la distance à zéro est le nombre lui-même.\par

\end{blocEdition}
\editionPage{indices-4}{Indications progressives}{Défis 4 à 6}{Ne lire que l’aide nécessaire.}
\begin{blocEdition}{indices-4}
\editionEtape{D04 · Remonter un carnet de bord}\textbf{Indice 1.} Calculer la somme des trois variations.\par\medskip\textbf{Indice 2.} Le départ est la valeur finale moins la variation totale.\par
\editionEtape{D05 · Une fraction entre deux autres}\textbf{Indice 1.} Écrire les deux bornes avec le dénominateur 15.\par\medskip\textbf{Indice 2.} Tester les numérateurs strictement entre 9 et 12.\par
\editionEtape{D06 · Un ruban en trois parts}\textbf{Indice 1.} Calculer la fraction restante avec le dénominateur 21.\par\medskip\textbf{Indice 2.} 32 cm correspondent à huit parts élémentaires.\par

\end{blocEdition}
\editionPage{indices-7}{Indications progressives}{Défis 7 à 9}{Ne lire que l’aide nécessaire.}
\begin{blocEdition}{indices-7}
\editionEtape{D07 · Le cadre autour d’un carré}\textbf{Indice 1.} Soustraire deux aires de carrés.\par\medskip\textbf{Indice 2.} Découper en quatre bandes de 8 cm sur 1 cm et quatre petits carrés de 1 cm de côté.\par
\editionEtape{D08 · Retrouver l’étape d’un motif}\textbf{Indice 1.} Retirer l’allumette supplémentaire.\par\medskip\textbf{Indice 2.} Le nombre restant doit pouvoir se partager en groupes de trois.\par
\editionEtape{D09 · La somme de cinq consécutifs}\textbf{Indice 1.} Écrire les nombres n, n+1, n+2, n+3, n+4.\par\medskip\textbf{Indice 2.} Réduire la somme puis chercher le facteur 5.\par

\end{blocEdition}
\editionPage{indices-10}{Indications progressives}{Défis 10 à 12}{Ne lire que l’aide nécessaire.}
\begin{blocEdition}{indices-10}
\editionEtape{D10 · Un partage dessiné}\textbf{Indice 1.} Retrouver la valeur d’une part.\par\medskip\textbf{Indice 2.} On peut aussi ajouter directement trois fois cinq au total initial.\par
\editionEtape{D11 · Deux réductions successives}\textbf{Indice 1.} La seconde réduction porte sur 80 euros.\par\medskip\textbf{Indice 2.} Comparer le prix final au prix initial.\par
\editionEtape{D12 · Un modèle et son domaine}\textbf{Indice 1.} Utiliser 24-3t.\par\medskip\textbf{Indice 2.} Une réserve ne peut pas avoir un volume négatif.\par

\end{blocEdition}
\editionPage{indices-13}{Indications progressives}{Défis 13 à 15}{Ne lire que l’aide nécessaire.}
\begin{blocEdition}{indices-13}
\editionEtape{D13 · Quatre angles et une bissectrice}\textbf{Indice 1.} Un tour vaut 360 degrés.\par\medskip\textbf{Indice 2.} Partager seulement l’angle concerné en deux.\par
\editionEtape{D14 · Une aire en plusieurs partages}\textbf{Indice 1.} Les triangles ont la même hauteur issue de C.\par\medskip\textbf{Indice 2.} Leurs bases sur (AB) valent 6, 3 et 3 cm.\par
\editionEtape{D15 · Un triangle et son centre}\textbf{Indice 1.} Quel est le milieu de [AA′] ? Et de [BB′] ?\par\medskip\textbf{Indice 2.} Ce sont les deux diagonales du quadrilatère indiqué.\par

\end{blocEdition}
\editionPage{indices-16}{Indications progressives}{Défis 16 à 18}{Ne lire que l’aide nécessaire.}
\begin{blocEdition}{indices-16}
\editionEtape{D16 · Même base, plusieurs aires}\textbf{Indice 1.} L’aire dépend de la base et de la hauteur.\par\medskip\textbf{Indice 2.} La longueur des côtés obliques peut changer.\par
\editionEtape{D17 · Deux cylindres de même capacité}\textbf{Indice 1.} Écrire le volume exact de A.\par\medskip\textbf{Indice 2.} L’aire de base de B est quatre fois celle de A.\par
\editionEtape{D18 · Une moyenne à atteindre}\textbf{Indice 1.} Une moyenne de 13 sur cinq notes exige une somme de 65.\par\medskip\textbf{Indice 2.} Pour maximiser la moyenne, prendre la cinquième note égale à 20.\par

\end{blocEdition}
\editionPage{indices-19}{Indications progressives}{Défis 19 à 20}{Ne lire que l’aide nécessaire.}
\begin{blocEdition}{indices-19}
\editionEtape{D19 · Deux conditions pour une urne}\textbf{Indice 1.} Trouver séparément les effectifs rouge et bleu.\par\medskip\textbf{Indice 2.} Les boules restantes seront vertes.\par
\editionEtape{D20 · Un escalier en boucle}\textbf{Indice 1.} Chaque répétition ajoute 30 à droite et 20 en haut.\par\medskip\textbf{Indice 2.} Le nombre de répétitions doit convenir aux deux déplacements.\par

\end{blocEdition}
